\scanpage{80}
% @1 first line and interline
3) Soit $u:A'\longrightarrow A$ une \deleted{isogénie}\inserted{\unsure{isom|hom|homom.}, avec $\dim A\unsure{.|-|=}\dim A'=n$}, alors
% @2–4 top-degree comparison
\diagram{\begin{tikzcd}[ampersand replacement=\&,column sep=small,row sep=small]
 u^{2n}: \& H^{2n}(A)\arrow[r]\arrow[d,phantom,"\rotatebox{90}{$\simeq$}" description] \& H^{2n}(A')\arrow[d,phantom,"\rotatebox{90}{$\simeq$}" description] \& \text{est donné par}\\
 \& \unsure{\mathbb Q|Q}_\ell(-n) \& \mathbb Q_\ell(-n) \&
\end{tikzcd}}
% @5–6 degree
la multiplication par $\deg u=d$ (nul ssi $u(A')\ne A$, i.e. $u$ n'est pas une isogénie).

% @7–8 transpose
En d'autres termes, définissant $u_*$
\[
 u_*:\ H^{*}(A)\mathrel{-}H^{\unsure{*|1}}(A')
\]
% @9–10 proviso
par transposition (Attention, à l'hyp que $A$ et $A'$ ont \unsure{=|même} dim) on trouve
% @11 unit
\[
 u_0(1_A)=d.1_{A'}
\]
% @12–13 arbitrary dimensions
Supposons que \deleted{les} \deleted{degrés} dim $n,n'$ de $A,A'$ soient quelconques, on définit par transposition
% @14 shifted map
\[
 H^i(A)\longrightarrow H^{i+2(n'-n)}(A')(n'-n)
\]
% @15–16 crossed opening
\deleted{et on va} qui se déduit \unsure{directmt|directement} de la connaissance de
% @17–18 pullback
\[
 u^1:\ H^1(A')\longrightarrow H^1(A)
 \qquad\left(\begin{gathered}\text{d'où }u^*:H^*(A')\to H^*(A)\\ \text{par }\bigwedge\nolimits^*\end{gathered}\right)
\]
% @19–20 last lines
et des isom $\unsure{\eta|\int}_A$ et $\eta_{A'}$, par transposition.

\unsure{Noter|Notons} que
% archive sheet number: 79

